\section{Related Work}
\label{sec:related}

\paragraph{Generative molecular docking.}
Diffusion-based docking models formulate ligand translation, rotation, and
torsion as a generative process on a non-Euclidean pose space
\citep{corso2023diffdock}. \method instead uses flow matching
\citep{lipman2023flow} and groups ligand atoms into rigid fragments, yielding a
product of fragment-level \(\SE(3)\) states.

\paragraph{Equivariant geometric learning.}
Equivariant graph networks represent geometric features using irreducible
representations and tensor products, allowing vector and higher-order features
to transform consistently under rigid motions
\citep{liao2023equiformer}. \method uses this machinery in a single
protein--ligand graph and predicts atom-level equivariant forces.

\paragraph{Docking datasets and validity.}
We train on PLINDER \citep{durairaj2024plinder} and evaluate frozen
reference-pocket diagnostics on Astex Diverse \citep{hartshorn2007astex},
PoseBusters v2 \citep{buttenschoen2024posebusters}, and CASF-2016
\citep{su2019casf}. We report pose accuracy separately from PoseBusters
physical-validity checks.

\todo{Expand this section only after selecting the final baseline set and
ensuring that every comparison uses a compatible pocket-definition protocol.}
